\documentclass[12pt,a4paper]{article}
\usepackage{ctex}
\usepackage{amsmath,amssymb}
\usepackage{graphicx}
\usepackage{geometry}
\usepackage{multirow}
\usepackage{caption}
\usepackage{enumitem}
\usepackage{fancyhdr}
\usepackage{lastpage}
\usepackage{float}

% 页面设置
\geometry{left=2.5cm,right=2.5cm,top=2.5cm,bottom=2.5cm}

% 页眉页脚
\pagestyle{fancy}
\fancyhf{}
\lhead{数字逻辑实验}
\rhead{实验一 组合逻辑电路}
\cfoot{第 \thepage 页 / 共 \pageref{LastPage} 页}

\begin{document}

\begin{center}
\LARGE \textbf{实验一 组合逻辑电路设计与仿真}
\end{center}

\vspace{1em}

\begin{tabular}{ll}
\textbf{实验名称：} & 组合逻辑电路设计与Multisim仿真 \\
\textbf{实验日期：} & \rule{5cm}{0.5pt} \\
\textbf{学生姓名：} & \rule{5cm}{0.5pt} \\
\textbf{学\qquad 号：} & \rule{5cm}{0.5pt} \\
\textbf{班\qquad 级：} & \rule{5cm}{0.5pt} \\
\textbf{指导教师：} & \rule{5cm}{0.5pt} \\
\end{tabular}

\vspace{1.5em}

\section{一、实验检查记录表}
\begin{table}[H]
\centering
\renewcommand{\arraystretch}{1.2}
\begin{tabular}{|c|c|c|c|c|}
\hline
检查项目 & 检查内容 & 自查结果 & 教师检查 & 备注 \\
\hline
实验预习 & 掌握组合逻辑设计流程、芯片原理 & & & \\
\hline
器件认知 & 识别74LS00、74LS138、74LS86等芯片 & & & \\
\hline
电路搭建 & 接线规范、电源接地连接正确 & & & \\
\hline
仿真操作 & 熟练使用Multisim完成电路仿真 & & & \\
\hline
数据测试 & 真值表数据记录完整准确 & & & \\
\hline
实验纪律 & 遵守实验要求，操作规范有序 & & & \\
\hline
实验报告 & 内容完整、书写条理清晰 & & & \\
\hline
\end{tabular}
\end{table}

\section{二、实验原理}
本实验研究组合逻辑电路，电路输出仅取决于当前输入信号，与过往状态无关，不具备记忆功能。

与非门是基础逻辑器件，逻辑表达式为 $Y=\overline{A\cdot B}$，逻辑特点为有0出1、全1出0，可组合搭建各类逻辑电路。实验采用74LS00二输入与非门芯片实现逻辑运算。

门电路存在传输延迟特性，平均传输延迟时间计算公式 $t_{pd}=\frac{1}{2}(t_{pHL}+t_{pLH})$，该参数用于衡量电路响应速度。

组合逻辑电路标准设计流程：明确设计要求、列出真值表、推导逻辑表达式、化简函数、选用器件搭建电路。

74LS138为3线-8线译码器，满足使能条件 $S_1=1,\overline{S_2}=\overline{S_3}=0$ 时正常译码，输出低电平有效，可利用最小项表达式实现三变量逻辑函数。

一位全加器完成一位二进制带进位加法运算，输入为加数、被加数、低位进位，输出为本位和与向高位进位，可借助译码器与门电路实现加法逻辑。

\section{三、实验器材}
\begin{itemize}[itemsep=2pt]
    \item 计算机一台
    \item Multisim 仿真软件
    \item 74LS00、74LS86、74LS138逻辑芯片
    \item LED指示灯、限流电阻、逻辑开关、电源与地
    \item 虚拟示波器
\end{itemize}

\section{四、实验内容与步骤}
\subsection{1. 裁判表决电路设计}
依据表决规则推导逻辑式 $Y=\overline{\overline{AB}\cdot\overline{AC}}$，使用74LS00搭建电路，切换输入电平观测LED输出状态，核对真值表。

\subsection{2. 译码器实现一位全加器}
将输入信号接入74LS138地址端，配置有效使能电平，通过与非门组合最小项，分别输出加法和值与进位信号，完成功能验证。

\subsection{3. 发电机控制组合电路设计}
以车间工作状态为输入，搭建异或门与与非门电路，按照开工数量对应发电机启停逻辑，仿真测试全部输入工况。

\section{五、实验结果与仿真波形}
各组电路仿真输出均与理论真值表保持一致，裁判电路、全加器、发电机控制电路均可按照预设逻辑正常工作，电平指示响应准确，电路设计符合功能要求。

\section{六、实验分析与总结}
\begin{enumerate}[itemsep=3pt]
    \item 组合逻辑电路输出仅由当前输入决定，无记忆存储特性。
    \item 译码器使能电平、芯片供电接地必须正确连接，电路方可正常工作。
    \item 逻辑函数化简能够有效减少门电路使用数量，简化硬件结构。
    \item Multisim仿真可快速排查接线与逻辑错误，提升实验效率。
    \item 负载类指示灯件需搭配限流电阻使用，避免器件损坏。
\end{enumerate}

\textbf{实验结论：}
本次完成多款典型组合逻辑电路的分析、设计与仿真测试，熟练掌握常用集成门电路与译码器件应用方法，熟悉组合电路完整设计流程，达成全部实验学习目标。

\vspace{2em}
\begin{flushright}
\textbf{实验人：}\rule{4cm}{0.5pt}\\
\textbf{日期：}\rule{4cm}{0.5pt}
\end{flushright}

\end{document}